% resume-entrymeta-inline.tex — smoke fixture: entrymeta=inline is a supported
% value and must build cleanly (issue #230). The entries cover the joins and the
% optional-field shapes that could put a separator next to an absent cell; what
% reaches the page is asserted by tests/extraction, what the option resolves to
% by tests/regression/components-entrymeta, and how it looks is a visual check.
\documentclass[fontsize=10pt, margin=normal, entrymeta=inline]{careerdossier-resume}
\CDossierSetup{
  name     = {Ada Lovelace},
  headline = {Data Scientist},
  email    = {ada@example.com},
}
\begin{document}
\MakeCDossierHeader
\CDossierSection{Experience}
\begin{CDossierEntry}[
  title        = {Senior Data Scientist},
  organization = {Example Corp},
  location     = {Berlin, DE},
  dates        = {2020--2024}
]
  \begin{CDossierItemize}
    \item Built a reproducible analytical workflow.
  \end{CDossierItemize}
\end{CDossierEntry}
\begin{CDossierEntry}[
  title    = {Volunteer Reviewer},
  location = {Remote}
]
  \begin{CDossierItemize}
    \item A location with no organization opens its line with no separator.
  \end{CDossierItemize}
\end{CDossierEntry}
\begin{CDossierEntry}[
  title        = {Workshop Facilitator},
  organization = {Example Trust}
]
  \begin{CDossierItemize}
    \item An organization with no location ends its line with no separator.
  \end{CDossierItemize}
\end{CDossierEntry}
\end{document}
